\section*{الفصل \ref{ch:g7:states-of-matter} --- حالات المادة وتغيّراتها}

\begin{solution}{exo:g7:states-of-matter:1}
الصلبة: \omterm{def:g7:states-of-matter:molecule}{جزيئات} مرصوصة في رصّة منظَّمة، ترتعش في مواضعها.
والسائلة: \omterm{def:g7:states-of-matter:molecule}{جزيئات} متلامسة غير منظَّمة، تنزلق بعضها بجوار بعض.
والغازية: \omterm{def:g7:states-of-matter:molecule}{جزيئات} متباعدة، تطير سريعة في كل الاتجاهات.
\end{solution}

\begin{solution}{exo:g7:states-of-matter:2}
في الحالة الصلبة يحفظ كل \omterm{def:g7:states-of-matter:molecule}{جزيء} موقعه في الرصّة، فيحفظ الكلّ
شكله. وفي السائلة تنزلق \omterm{def:g7:states-of-matter:molecule}{الجزيئات} بحرّية بعضها بجوار بعض وهي
متلامسة: فينهار الحشد ليلائم الإناء، حتى لا يستطيع شيء أن
ينزلق أخفض --- وهو السطح المستوي.
\end{solution}

\begin{solution}{exo:g7:states-of-matter:3}
\omterm{def:g2:air-around-us:air}{الهواء} أكثره حيّز خالٍ بين \omterm{def:g7:states-of-matter:molecule}{جزيئات} طائرة: فالدفع يرصّ الطائرين
أقرب (وطرقها يقاوم --- وهو المرونة). أما \omterm{def:g7:states-of-matter:molecule}{جزيئات} الماء فمتلامسة
أصلًا؛ ولا خواء يُعصر، فيلقى المكبس جدارًا.
\end{solution}

\begin{solution}{exo:g7:states-of-matter:4}
يثبت أن \omterm{def:g7:states-of-matter:molecule}{جزيئات} \omterm{def:g3:solids-liquids-gases:gas}{الغاز} تطير متباعدة وبينها خواء عظيم --- فيمكن
رصّ ملء غرف معًا. أما الماء، وجزيئاته متلامسة أصلًا، فغير \omterm{prop:g7:states-of-matter:squeeze}{قابل
للانضغاط}: فلا منفاخ يرصّه أضيق.
\end{solution}

\begin{solution}{exo:g7:states-of-matter:5}
الدفء يجعل \omterm{def:g7:states-of-matter:molecule}{الجزيئات} المرصوصة ترتعش أشدّ فأشدّ؛ وعند \omterm{def:g4:changes-of-state:melting-point}{درجة
الانصهار} يغلب الارتعاش قبضة الرصّة، فتنكسر الصفوف، وتبدأ
\omterm{def:g7:states-of-matter:molecule}{الجزيئات} بالانزلاق: فيصير الصلب سائلًا.
\end{solution}

\begin{solution}{exo:g7:states-of-matter:6}
تغيّر الحالة لا يفعل إلا أن يعيد ترتيب \omterm{def:g7:states-of-matter:molecule}{الجزيئات} نفسها --- فلا
يُخلق منها شيء ولا يُفنى، فلا تستطيع \omterm{def:g3:measuring-mass:mass}{الكتلة} أن تتغيّر. وفقدان
\omterm{def:g3:measuring-mass:mass}{الكتلة} كان يقتضي اختفاء \omterm{def:g7:states-of-matter:molecule}{جزيئات} (أو هروبها من الإناء --- وهو
إنذار القنّينة المفتوحة الكاذب).
\end{solution}

\begin{solution}{exo:g7:states-of-matter:7}
\omterm{def:g6:measurement-in-science:measuring}{تقيس} نشاط الحركة الجزيئية. وحين يلمس الساخنُ الباردَ تقصف
\omterm{def:g7:states-of-matter:molecule}{الجزيئات} النشيطة الخاملةَ، فتسلّمها النشاط تصادمًا تصادمًا ---
فتهدأ النشيطة وتنشط الخاملة --- حتى يتقاسم الحشدان وتيرة
واحدة: أي تساوي درجتَي \omterm{def:g5:heat-and-insulation:heat}{الحرارة}.
\end{solution}

\begin{solution}{exo:g7:states-of-matter:8}
\omterm{def:g7:states-of-matter:molecule}{جزيئات} العطر تطير، وتتصادم \omterm{def:g7:states-of-matter:molecule}{بجزيئات} \omterm{def:g2:air-around-us:air}{الهواء}، وترتدّ، وتهيم ---
زيغزاغًا بعد زيغزاغ عبر الغرفة. والنسيم يحمل حشودًا كاملة منها
حملًا: فركوب \omterm{def:g2:air-around-us:wind}{الريح} خير من التعثّر في الزحام.
\end{solution}

\begin{solution}{exo:g7:states-of-matter:9}
لا تستطيع أن تنتزع نفسها من قبضة \omterm{def:g3:solids-liquids-gases:liquid}{السائل} إلا \emph{أسرع}
\omterm{def:g7:states-of-matter:molecule}{الجزيئات}، \omterm{def:g6:water-states:changes}{فالتبخّر} ضريبة على النشاط: إذ يحمل الهاربون أكثر من
حصّتهم، ويبقى الحشد وراءهم أبطأ وسطيًّا --- أي أبرد. فالعرق
والإصبع المبلولة يبرّدان بتصدير أنشط ما عندهما.
\end{solution}

\begin{solution}{exo:g7:states-of-matter:10}
في الشمس تطير \omterm{def:g7:states-of-matter:molecule}{الجزيئات} المحبوسة أسرع وتطرق الغلاف أشدّ وأكثر:
فتنفخ عاصفة البرَد البالونَ حتى يوازنها الغلاف المشدود. وفي
البرد تبطئ \omterm{def:g7:states-of-matter:molecule}{الجزيئات} نفسها، فيضعف الطرق، ويهرس دفعُ \omterm{def:g2:air-around-us:air}{الهواء}
الخارجي البالونَ أصغر.
\end{solution}

\begin{solution}{exo:g7:states-of-matter:11}
الماء المتجمّد ينغلق في رصّة هوائية مفتوحة غير معتادة --- فيها
حيّز خالٍ أكثر ممّا كان في الحشد المنزلق --- فيشغل الثلج حيّزًا
أكبر \omterm{def:g7:states-of-matter:molecule}{للجزيئات} نفسها: فكثافته دون \omterm{def:g6:density-floating:density}{كثافة} \omterm{def:g3:solids-liquids-gases:liquid}{السائل}، \omterm{def:g1:floating-and-sinking:words}{فيطفو}. ولو كان
الماء عاديًّا لرصّ الثلج ضيّقًا، ولغرق وهو يتكوّن، ولتجمّدت
البِرَك من القاع إلى الأعلى --- ولم يبقَ للأسماك ماء محميّ.
\end{solution}

\begin{solution}{exo:g7:states-of-matter:12}
لم يحرّك الماءَ شيء، ومع ذلك سافر الحبر: \omterm{def:g7:states-of-matter:molecule}{فجزيئات} \omterm{def:g3:solids-liquids-gases:liquid}{السائل} نفسه في
حركة لا تنقطع، تزاحم \omterm{def:g7:states-of-matter:molecule}{جزيئات} الحبر خطوةً عشوائية بعد خطوة عبر
الحشد --- فكلا السائلين حيّ بالحركة في كل لحظة، ولو بدت الكأس
ساكنة سكونًا تامًّا. والدفء يسرّع كل مزاحمة، فيتعجّل الانتشار.
وهكذا جعل الحبر الصبور الرقصةَ الخفيّة مرئية --- دليلًا أقنع
المشكّكين في \omterm{def:g7:states-of-matter:molecule}{الجزيئات} نفسها.
\end{solution}

\begin{solution}{pb:g7:states-of-matter:1}
\textbf{1.} في صفوف مرتّبة، كتفًا إلى كتف، وكل واحد يهتزّ في
موضعه --- ولا يجوز لأحد أن يبادل موضعه أو يغادر موقعه.
\textbf{2.} تذوب الصفوف: فيبقى التلاميذ كتفًا إلى كتف
(متلامسين!) لكنهم ينسجون وينزلقون بعضهم بجوار بعض. وعليهم أن
يحفظوا التماسّ مع الحشد --- فلا يجوز لأحد أن يجري طليقًا.
\textbf{3.} يتبعثر التلاميذ، ويجرون سريعًا مستقيمًا حتى
يتصادموا أو يصطدموا بجدار، فيرتدّون بلا نهاية --- وإذ لا شيء
يمسكهم معًا فلا يوقفهم إلا الجدران: فيطالب الحشد بكل ركن
تعرضه الصالة.
\textbf{4.} نعم --- الثلاثون أنفسهم في كل مشهد: \omterm{def:g3:measuring-mass:mass}{فالكتلة} تنجو
من تغيّرات الحالة، ممثَّلةً.
\textbf{5.} يسمح مشهد \omterm{def:g3:solids-liquids-gases:gas}{الغاز} بذلك --- فيتزاحم الراكضون في نصف
الصالة (ويطرقون السائقين أشدّ). أما كومة \omterm{def:g3:solids-liquids-gases:liquid}{السائل} فلا تستطيع أن
تنكمش: فالتلاميذ متلامسون أصلًا. والممثَّل: الغازات تنضغط
والسوائل لا.
\textbf{6.} أن يهتزّوا في مواقعهم أشدّ فأشدّ؛ ويصير المشهد
سائلًا في اللحظة التي يفجّر فيها \omterm{def:g3:sound-around-us:vibration}{الاهتزاز} الصفوف ويبدأ
التلاميذ بالانزلاق بعضهم بجوار بعض.
\textbf{7.} تُنفَق \omterm{def:g5:everyday-energy:energy}{طاقة} القرص على كسر قبضة الرصّة --- صفًّا
صفًّا --- لا على حركة أسرع: فعلى الناقد \omterm{def:g3:temperature-thermometers:temperature}{ميزان الحرارة} أن يبلّغ
بأن النشاط الوسطيّ \emph{لم يتغيّر} حتى ينكسر آخر صفّ. وهي
المسطبة، ممثَّلةً.
\textbf{8.} \emph{أسرع} التلاميذ --- فهم وحدهم من يستطيع
الانتزاع. وينخفض نشاط الكومة الباقية الوسطيّ: \omterm{def:g3:solids-liquids-gases:liquid}{فالسائل} يبرّد
نفسه \omterm{def:g6:water-states:changes}{بالتبخّر}، وهو بعينه قشعريرة العرق والأصابع المبلولة.
\textbf{9.} يواصل الممثّلون التصادم بالحبل ودفعه إلى الخارج؛
والممثّلون الأسرع يدفعون أشدّ وأكثر، فتنتفخ الحلقة --- وهو
البالون في الشمس.
\textbf{10.} ``ارصّوا على نظام \emph{مفتوح} --- بمسافة ذراع،
هوائية'': فصلب الماء يشغل حيّزًا أكبر من سائله، \omterm{def:g1:floating-and-sinking:words}{فيطفو} طوف مشهد
الثلج على حشد مشهد \omterm{def:g3:solids-liquids-gases:liquid}{السائل}.
\textbf{11.} يُزاحَم القميص الأزرق --- ويُدفَع خطوةً عشوائية
بعد خطوة عبر الحشد المنزلق حتى يكون قد هام في كل مكان: وهو
انتشار الحبر. وارفع القرص تتعجّل كل مزاحمة، فيتعجّل الهيام.
\textbf{12.} مثلًا: ``مسرحية الليلة تمثّل فكرة واحدة: أن كل
مادة \omterm{def:g7:states-of-matter:molecule}{جزيئات} في حركة. وبالتسليم بها يأتي قانونان قديمان
بالمجّان --- نجاة \omterm{def:g3:measuring-mass:mass}{الكتلة} من كل تغيّر حالة، ومسطبة \omterm{def:g6:water-states:changes}{الانصهار}.
وما لا نستطيع تمثيله هو المقياس: فالممثّلون الحقيقيون أصغر من
كل رؤية وأكثر من جماهير العالم كلها.''
\end{solution}
